% Part: propositional-logic
% Chapter: syntax-and-semantics
% Section: introduction

\documentclass[../../../include/open-logic-section]{subfiles}

\begin{document}

\olfileid{pl}{syn}{int}

\olsection{పరిచయం}

ప్రతిజ్ఞావాక్యాత్మక తర్కం, \tetoken{ప్రతిజ్ఞావాక్య చరాలు}{!!{propositional variable}s} నుంచి $\lnot$, $\land$,
$\lor$, $\lif$, $\liff$ అనే ప్రతిజ్ఞావాక్య సంయోజకాలను ఉపయోగించి నిర్మించిన
\tetoken{సూత్రాలతో}{!!{formula}s} వ్యవహరిస్తుంది. అంతర్బుద్ధి ప్రకారం,
\tetoken{ఒక ప్రతిజ్ఞావాక్య చరం}{!!a{propositional variable}}~$p$ అనేది సత్యమో అసత్యమో అయ్యే ఒక వాక్యం
లేదా ప్రతిజ్ఞావాక్యానికి ప్రతినిధిగా ఉంటుంది. ఒక \tetoken{సూత్రంలోని}{!!a{formula}}
\tetoken{ప్రతిజ్ఞావాక్య చరం}{!!{propositional variable}} యొక్క “సత్యమూల్యం” నిర్ణయమైనప్పుడల్లా,
వాటిని ప్రతిజ్ఞావాక్య సంయోజకాలతో కలిపి ఏర్పరిచిన ఏ \tetoken{సూత్రాలు}{!!{formula}s} సత్యమూల్యమైనా
నిర్ణయమవుతుంది. దాని అర్థవిచారం సత్యమూల్యాల ప్రమేయాల ద్వారా
ఇవ్వబడుతుంది కనుక ప్రతిజ్ఞావాక్యాత్మక తర్కాన్ని \emph{సత్యమూల్య-ప్రమేయాత్మకం}
అంటాం. ముఖ్యంగా, ప్రతిజ్ఞావాక్యాత్మక తర్కంలో సత్యాసత్యాల గురించి మరింత సున్నితమైన
నిర్ణయాలను పక్కన పెడతాం: ఉదాహరణకు, ఒక విషయం కేవలం
పరిస్థిత్యాధీనంగా కాకుండా అనివార్యంగా సత్యమా, అది సత్యమని మనకు
తెలుసా, లేదా అది గతంలో సత్యమా, ఇప్పుడు సత్యమా, భవిష్యత్తులో సత్యం
అవుతుందా అనే భేదాలను పరిగణించం. సత్యం ($\True$), అసత్యం
($\False$) అనే రెండు సత్యమూల్యాలను మాత్రమే పరిగణిస్తాం; అందువల్ల ఒక
ప్రకటన సత్యమూ అసత్యమూ కాకపోవచ్చు లేదా సగం మాత్రమే సత్యం కావచ్చు అనే
అవకాశాన్ని చర్చ నుంచి తొలగిస్తాం. వాటితో నిర్మించిన \tetoken{ఒక సూత్రం}{!!a{formula}}
సత్యమూల్యం దాని భాగాల సత్యమూల్యాల ద్వారానే పూర్తిగా నిర్ణయమయ్యే
సంయోజకాలపైనే దృష్టి పెడతాం; ఉదాహరణకు, దాని అర్థంపై ఆధారపడే
సంయోజకాలను ఇక్కడ పరిగణించం. ముఖ్యంగా, ఆంగ్లంలోని షరతీయ వాక్యాల
సత్యమూల్యం ఈ అర్థంలో సత్యమూల్య-ప్రమేయాత్మకమా అన్నది వివాదాస్పదం.
భౌతిక షరతీయ సంయోజకం~$\lif$ మాత్రం అలాంటిదే; సత్యమూల్య-ప్రమేయాత్మకం
కాని షరతీయ సంయోజకాలతో ఇతర తర్కవ్యవస్థలు వ్యవహరిస్తాయి.

సత్యమూల్య-ప్రమేయాత్మక ప్రతిజ్ఞావాక్యాత్మక తర్క సిద్ధాంతాన్ని, అధిసిద్ధాంతాన్ని
అభివృద్ధి చేయాలంటే, ముందుగా దాని వ్యక్తీకరణల వాక్యనిర్మాణం,
అర్థవిచారాన్ని నిర్వచించాలి. సంయోజకాలను ఉపయోగించి
\tetoken{ప్రతిజ్ఞావాక్య చరాలు}{!!{propositional variable}s} నుంచి \tetoken{సూత్రాలను}{!!{formula}s} నిర్మించే ఒక
విధానాన్ని వివరిస్తాం. ప్రత్యామ్నాయ నిర్వచనాలు కూడా సాధ్యమే. ఇతర
వ్యవస్థలు వేరే సంకేతాలను ఎంచుకోవచ్చు, వేరే సంయోజకాల సమితులను
ప్రాథమికమైనవిగా తీసుకోవచ్చు, కుండలీకరణలను వేరుగా వాడవచ్చు (లేదా
పోలిష్ సంజ్ఞావిధానంలోలాగా అసలు వాడకపోవచ్చు). అయితే అన్ని విధానాల్లో
ఉమ్మడి విషయం ఏమిటంటే, నిర్మాణ నియమాలు \tetoken{సూత్రాలు}{!!{formula}s} సమితిని
\emph{ఆగమనాత్మకంగా} నిర్వచిస్తాయి. సరిగ్గా చేస్తే, ప్రతి వ్యక్తీకరణ
నిర్మాణ నియమాల ప్రకారం సారాంశంగా ఒకే విధంగా ఏర్పడుతుంది. ఫలిత
వ్యక్తీకరణలు \emph{ఏకైకంగా చదవదగినవి} అయ్యే ఈ ఆగమనాత్మక నిర్వచనం
వల్ల, అదే పద్ధతితో---ఆగమనాత్మక నిర్వచనంతో---వాటికి అర్థాలు
ఇవ్వగలం.

వ్యక్తీకరణలకు అర్థం ఇవ్వడం అర్థవిచార పరిధి. ప్రతిజ్ఞావాక్యాత్మక తర్క
అర్థవిచారంలో కేంద్ర భావన, ఒక \tetoken{సత్యమూల్య కేటాయింపులో}{!!a{valuation}} సంతృప్తి.
\tetoken{ఒక సత్యమూల్య కేటాయింపు}{!!^a{valuation}}~$\pAssign{v}$, \tetoken{ప్రతిజ్ఞావాక్య చరాలకు}{!!{propositional variable}s}
$\True$, $\False$ అనే సత్యమూల్యాలను కేటాయిస్తుంది. ఏ
\tetoken{సత్యమూల్య కేటాయింపు}{!!{valuation}} అయినా ప్రతి \tetoken{సూత్రం}{!!{formula}}~$!A$కు
$\pValue{v}(!A)$ అనే సత్యమూల్యాన్ని నిర్ణయిస్తుంది. ఒక
\tetoken{సూత్రం}{!!^a{formula}}, \tetoken{ఒక సత్యమూల్య కేటాయింపు}{!!a{valuation}}~$\pAssign{v}$లో సంతృప్తమవడం
$\pValue{v}(!A) = \True$ అయినప్పుడూ, అప్పుడు మాత్రమే; దీనిని
$\pSat{v}{!A}$గా రాస్తాం. ఈ
సంబంధాన్ని కూడా $!A$ నిర్మాణంపై ఆగమనం ద్వారా నిర్వచించవచ్చు:
తార్కిక సంయోజకాల సత్య ప్రమేయాలను ఉపయోగించి, ఉదాహరణకు
$!A \land !B$ సంతృప్తిని $!A$, $!B$ల సంతృప్తి (లేదా
అసంతృప్తి) ద్వారా నిర్వచించవచ్చు.

వాక్యాల కోసం $\pSat{v}{!A}$ అనే సంతృప్తి సంబంధం ఆధారంగా
సర్వసత్యం, అర్థపర అనుగమనం, సంతృప్తిపరచదగినతనం అనే ప్రాథమిక అర్థపర
భావనలను నిర్వచించవచ్చు. ప్రతి \tetoken{సత్యమూల్య కేటాయింపు}{!!{valuation}} దానిని సంతృప్తిపరిస్తే,
అంటే ఏ~$\pAssign{v}$కైనా $\pValue{v}(!A) = \True$ అయితే,
\tetoken{ఒక సూత్రం}{!!^a{formula}} సర్వసత్యం, $\Entails !A$. $\Gamma$లోని అన్ని
\tetoken{సూత్రాలను}{!!{formula}s} సంతృప్తిపరిచే ప్రతి \tetoken{సత్యమూల్య కేటాయింపు}{!!{valuation}}, $!A$ను కూడా
సంతృప్తిపరిస్తే, ఆ సూత్రం ఆ \tetoken{సూత్రాలు}{!!{formula}s} సమితి నుంచి అర్థపరంగా అనుగమిస్తుంది:
$\Gamma \Entails !A$. ఏదో ఒక \tetoken{సత్యమూల్య కేటాయింపు}{!!{valuation}}, ఒక \tetoken{సూత్రాలు}{!!{formula}s}
సమితిలోని అన్ని \tetoken{సూత్రాలను}{!!{formula}s} ఒకేసారి సంతృప్తిపరిస్తే ఆ సమితి
సంతృప్తిపరచదగినది. \tetoken{సూత్రాలు}{!!{formula}s} ఆగమనాత్మకంగా నిర్వచించబడ్డాయి,
\tetoken{సూత్రాలు}{!!{formula}s} నిర్మాణంపై ఆగమనం ద్వారా సంతృప్తి కూడా నిర్వచించబడింది కాబట్టి,
మన అర్థవిచార ధర్మాలను నిరూపించడానికి, నిర్వచించిన అర్థపర భావనలను
పరస్పరం అనుసంధానించడానికి ఆగమనాన్ని ఉపయోగించవచ్చు.


\end{document}
